\chapter{ĐẶC TẢ TEST CASE}
\label{app:test-cases}

Phụ lục này định nghĩa đầy đủ các test case được nhắc đến trong Chương~6. Mỗi test case mô tả điều kiện đầu vào, các bước thực hiện và tiêu chí chấp nhận tương ứng.

\renewcommand{\arraystretch}{1.3}

% -------------------------------------------------------
\section*{TC-GRADE-01: Chấm phiếu chuẩn}
\addcontentsline{toc}{section}{TC-GRADE-01: Chấm phiếu chuẩn}

\begin{longtable}{|L{3.5cm}|L{10.5cm}|}
\hline
\textbf{Mã test case} & TC-GRADE-01 \\
\hline
\textbf{Tên} & Chấm một phiếu trắc nghiệm chuẩn \\
\hline
\textbf{Mục tiêu} & Xác nhận pipeline 15 bước xử lý đúng phiếu chụp trong điều kiện ánh sáng tốt, bốn marker rõ ràng và bong bóng tô đậm. \\
\hline
\textbf{Điều kiện tiên quyết} & Bài thi đã được tạo và có ít nhất một bộ đáp án; Form Profile tương ứng đã được cấu hình. \\
\hline
\textbf{Dữ liệu đầu vào} & Ảnh JPEG phiếu 40 câu, chụp thẳng góc, ánh sáng đồng đều, bốn marker đen góc rõ nét, học sinh tô đậm đáp án. \\
\hline
\textbf{Các bước thực hiện} &
\begin{enumerate}[leftmargin=1em]
  \item Đăng nhập, mở bài thi đã có đáp án.
  \item Chọn ảnh phiếu từ thiết bị và nhấn \textbf{Chấm bài}.
  \item Frontend gửi \path{POST /api/omr/grade} kèm ảnh và tham số bài thi.
  \item Backend chạy pipeline 15 bước, trả JSON kết quả và URL overlay.
  \item Frontend hiển thị điểm, MSSV, mã đề và ảnh overlay.
\end{enumerate} \\
\hline
\textbf{Kết quả kỳ vọng} &
\begin{itemize}[leftmargin=1em]
  \item HTTP 200, trường \texttt{success = true}.
  \item \texttt{uncertain\_count = 0} hoặc rất nhỏ (< 2 câu/40 câu).
  \item \texttt{warning\_codes} không chứa \texttt{MCQ\_MAP\_SEARCH\_RESCUE}.
  \item Ảnh overlay hiển thị ô đúng màu xanh, ô sai màu đỏ.
  \item Bản ghi lưu vào \path{omr\_grade\_result}, \path{last\_result} cập nhật.
\end{itemize} \\
\hline
\textbf{Test case liên quan} & TC-GRADE-02 \\
\hline
\end{longtable}

% -------------------------------------------------------
\section*{TC-GRADE-02: Chấm phiếu có câu không rõ (kích hoạt Rescue)}
\addcontentsline{toc}{section}{TC-GRADE-02: Chấm phiếu có câu không rõ}

\begin{longtable}{|L{3.5cm}|L{10.5cm}|}
\hline
\textbf{Mã test case} & TC-GRADE-02 \\
\hline
\textbf{Tên} & Chấm phiếu có bong bóng tô mờ và kích hoạt MCQ Map Search Rescue \\
\hline
\textbf{Mục tiêu} & Xác nhận cơ chế Rescue kích hoạt khi \texttt{uncertain\_count} vượt ngưỡng và cải thiện kết quả decode MCQ. \\
\hline
\textbf{Điều kiện tiên quyết} & Như TC-GRADE-01. Profile không có \texttt{disable\_mcq\_rescue = true}. \\
\hline
\textbf{Dữ liệu đầu vào} & Ảnh phiếu 40 câu chụp dưới ánh đèn huỳnh quang lệch, một số bong bóng tô mờ, lưới MCQ có thể lệch vài pixel sau warp. \\
\hline
\textbf{Các bước thực hiện} &
\begin{enumerate}[leftmargin=1em]
  \item Thực hiện như TC-GRADE-01 với ảnh phiếu có điều kiện khó.
  \item Quan sát trường \texttt{warning\_codes} và \texttt{uncertain\_count} trong JSON trả về.
  \item Kiểm tra file \path{bubble\_confidence\_json} trong \path{storage/uploads/omr/}.
\end{enumerate} \\
\hline
\textbf{Kết quả kỳ vọng} &
\begin{itemize}[leftmargin=1em]
  \item \texttt{uncertain\_count} baseline vượt $\max(4,\ \lfloor 0{,}60 \times 40 \rceil) = 24$.
  \item \texttt{warning\_codes} chứa \texttt{MCQ\_MAP\_SEARCH\_RESCUE}.
  \item Metadata ghi \texttt{line\_h\_before}, \texttt{line\_h\_after}, \texttt{top\_shift\_px}, \texttt{tested} (số ứng viên đã thử).
  \item \texttt{uncertain\_count} sau Rescue nhỏ hơn baseline ít nhất 2 câu (nếu Rescue thành công).
  \item Nếu không có ứng viên nào tốt hơn, kết quả baseline được giữ nguyên.
\end{itemize} \\
\hline
\textbf{Test case liên quan} & TC-GRADE-01 \\
\hline
\end{longtable}

% -------------------------------------------------------
\section*{TC-CAM-01: Smart Camera Scanner}
\addcontentsline{toc}{section}{TC-CAM-01: Smart Camera Scanner}

\begin{longtable}{|L{3.5cm}|L{10.5cm}|}
\hline
\textbf{Mã test case} & TC-CAM-01 \\
\hline
\textbf{Tên} & Tự động phát hiện phiếu và chụp ảnh bằng Smart Camera Scanner \\
\hline
\textbf{Mục tiêu} & Xác nhận máy trạng thái camera hoạt động đúng, chuyển từ \texttt{DETECTING} sang \texttt{LOCKED} sau bốn frame liên tiếp và tự chụp sau 1300ms. \\
\hline
\textbf{Điều kiện tiên quyết} & Trình duyệt chạy trên HTTPS hoặc localhost; thiết bị có camera; bài thi đã có Form Profile với \texttt{corner\_markers} được cấu hình. \\
\hline
\textbf{Dữ liệu đầu vào} & Luồng video camera hướng vào phiếu trắc nghiệm đặt trên mặt phẳng đủ sáng, bốn marker đen góc phiếu nằm trong khung viewfinder. \\
\hline
\textbf{Các bước thực hiện} &
\begin{enumerate}[leftmargin=1em]
  \item Mở tab \textbf{Camera} trong màn hình chấm bài.
  \item Nhấn \textbf{Bật camera} -- trạng thái chuyển sang \texttt{ACTIVE}.
  \item Hướng camera vào phiếu -- hệ thống phân tích frame mỗi 130ms (\texttt{DETECTING}).
  \item Căn phiếu vào khung: khi bốn \texttt{darkRatio} $\geq 0.14$ và \texttt{centerLuma} $\geq 52$ qua đủ bốn frame liên tiếp, khung chuyển xanh (\texttt{LOCKED}).
  \item Sau 1300ms ở trạng thái \texttt{LOCKED}, hệ thống tự gọi \texttt{canvas.toBlob()} -- ảnh thêm vào \texttt{pickedFiles} và trạng thái về \texttt{DETECTING}.
\end{enumerate} \\
\hline
\textbf{Kết quả kỳ vọng} &
\begin{itemize}[leftmargin=1em]
  \item Luồng video hiển thị đúng trong viewfinder.
  \item Khung chuyển xanh (class \texttt{match}) khi đủ điều kiện marker.
  \item Ảnh JPEG (quality 0.92) được thêm vào danh sách -- giáo viên không cần nhấn nút chụp.
  \item Khi môi trường không phải HTTPS/localhost, giao diện hiển thị cảnh báo rõ ràng thay vì lỗi im lặng.
\end{itemize} \\
\hline
\textbf{Test case liên quan} & TC-GRADE-01 (ảnh từ camera được xử lý bởi cùng pipeline) \\
\hline
\end{longtable}
